\documentclass[11pt, a4paper]{article}
\usepackage[a4paper, top=2.5cm, bottom=2.5cm, left=2cm, right=2cm]{geometry}
\usepackage{amsmath}
\usepackage{amssymb}

\title{\textbf{AXON-QIM: Quantum Information Manifold}\\ \large A Holographic Framework for Stochastic Resolution in Sub-2nm Lithography}
\author{AXON Research Group}
\date{\today}

\begin{document}
\maketitle

\begin{abstract}
As semiconductor manufacturing enters the High-NA EUV era ($<$2nm), traditional deterministic lithography models fail due to the "Stochastic Gap"---the dominant influence of photon shot noise. This paper derives the Quantum Information Manifold (AXON-QIM), a mathematical framework treating silicon resist as a Quantum Measurement Surface. By applying Holographic Principle and Von Neumann Entropy Minimization, we prove that stochastic defects are "Information Deficits" correctable via Field Inversion.
\end{abstract}

\section{The Information-Action Principle}
The "Action" $S[\phi]$ that governs the formation of a circuit line is:
\begin{equation}
    \delta S_{total} = \delta (S_{Quantum} + \beta S_{Info}) = 0
\end{equation}

\section{Derivation of the Information Force}
To recover the "Platonic Ideal", we apply gradient descent to the entropy manifold:
\begin{equation}
    F_{info} = -\nabla_{\phi} \left[ - \text{Tr}(\rho \ln \rho) \right]
\end{equation}

\section{The Inversion Operator}
The core AXON-QIM algorithm relies on the Stochastic Inversion Operator $\hat{\Omega}$:
\begin{equation}
    \hat{\Omega} = (\mathbf{I} - \mathbf{\Sigma}_{stoch})^{-1} \cdot \mathbf{K}_{holographic}
\end{equation}

\end{document}
